\documentclass[Journal,letterpaper]{ascelike}
\usepackage{graphicx}
\usepackage{booktabs}
\usepackage{amsmath}

\title{Title of the Paper}

\author[1]{First Author}
\author[2]{Second Author}
\affil[1]{Department of Civil Engineering, University, City, Country. Email: first@example.org}
\affil[2]{Company, City, Country}

\begin{document}

\maketitle

\begin{abstract}
State the problem, what you did, and the main result in a few sentences.
\end{abstract}

\KeyWords{first keyword, second keyword, third keyword}

\section{Introduction}
\label{sec:intro}
Introduce the problem and why it matters, then the contributions of this
paper. Cite related work like this~\cite{timoshenko1959}, or several works at
once~\cite{timoshenko1959,zienkiewicz2005}.

\section{Method}
\label{sec:method}
Describe what you did in enough detail to reproduce it. Numbered equations use
the \texttt{equation} environment:
\begin{equation}
  \label{eq:main}
  \sigma = \frac{M y}{I}.
\end{equation}
Refer to it as Equation~\eqref{eq:main}.

\section{Results}
\label{sec:results}
Table~\ref{tab:results} and Figure~\ref{fig:overview} show the main results.

\begin{table}[htbp]
  \centering
  \caption{Results on the test set.}
  \label{tab:results}
  \begin{tabular}{lrr}
    \toprule
    Method & Accuracy (\%) & Time (s) \\
    \midrule
    Baseline & 81.2 & 12 \\
    Ours & \textbf{87.5} & 15 \\
    \bottomrule
  \end{tabular}
\end{table}

\begin{figure}[htbp]
  \centering
  \includegraphics[width=0.6\linewidth]{example-image}
  \caption{Overview of the method. Upload an image and use its file name here.}
  \label{fig:overview}
\end{figure}

\section{Discussion}
Interpret the results, compare them with earlier work and state the limitations.

\section{Conclusion}
Summarise the findings and what comes next.

\section{Data Availability Statement}
Say which data, models or code support the findings and how to get them.

\section*{Acknowledgments}
Thank the people and funding that supported this work.

\bibliography{references}

\end{document}
